% Nini, mode d’emploi — compiler avec XeLaTeX (polices Inter et DejaVu installées) :
%   xelatex nini-mode-emploi.tex
\documentclass[11pt,a4paper]{article}

\usepackage[french]{babel}
\usepackage{fontspec}
\setmainfont{Inter}[
  UprightFont = *-Regular, BoldFont = *-SemiBold,
  ItalicFont = *-Italic, BoldItalicFont = *-SemiBoldItalic,
  Extension = .otf ]
\newfontfamily\displayfont{InterDisplay}[
  UprightFont = *-Bold, BoldFont = *-Bold, Extension = .otf ]
\setmonofont{DejaVu Sans Mono}[Scale=0.85]
\newfontfamily\symb{DejaVu Sans}

\usepackage[a4paper, margin=22mm, top=20mm, bottom=24mm]{geometry}
\usepackage{xcolor}
\usepackage{microtype}
\usepackage{enumitem}
\usepackage{tabularx}
\usepackage{colortbl}
\usepackage{titlesec}
\usepackage{needspace}
\usepackage{fancyhdr}
\usepackage{tikz}
\usetikzlibrary{babel}
\usepackage[most]{tcolorbox}
\usepackage[hidelinks, pdftitle={Nini, mode d'emploi}, pdfauthor={Nini}]{hyperref}

% ---------- Couleurs (celles de l'app) ----------
\definecolor{accent}{HTML}{3355FF}
\definecolor{accentsoft}{HTML}{EEF1FF}
\definecolor{ink}{HTML}{17171C}
\definecolor{ink2}{HTML}{52525C}
\definecolor{line}{HTML}{E2E3E9}
\definecolor{soft}{HTML}{F5F6F9}
\definecolor{ok}{HTML}{1D9A61}
\definecolor{oksoft}{HTML}{E9F7F0}
\definecolor{no}{HTML}{D93F45}
\definecolor{nosoft}{HTML}{FDEEEE}
\definecolor{amber}{HTML}{B8780A}
\definecolor{ambersoft}{HTML}{FFF5E3}
\definecolor{pink}{HTML}{E0569A}
\color{ink}

% ---------- Mise en page ----------
\setlength{\parindent}{0pt}
\setlength{\parskip}{6pt}
\linespread{1.15}
\titleformat{\section}{\needspace{9\baselineskip}\displayfont\fontsize{17}{21}\selectfont}{}{0pt}{}
\titlespacing*{\section}{0pt}{20pt}{8pt}
\pagestyle{fancy}
\fancyhf{}
\renewcommand{\headrulewidth}{0pt}
\fancyfoot[L]{\small\color{ink2}Nini · mode d'emploi}
\fancyfoot[R]{\small\color{ink2}\thepage}

% ---------- Petits composants ----------
\newcommand{\muted}[1]{{\color{ink2}#1}}
\newcommand{\kbd}[1]{\tcbox[on line, boxsep=0pt, left=3pt, right=3pt, top=1.5pt, bottom=1.5pt,
  colback=soft, colframe=line, boxrule=0.5pt, arc=2pt]{\ttfamily\small #1}}
\newcommand{\btn}[1]{\textbf{#1}}

% Étapes numérotées dans un rond bleu
\newcommand{\pastille}[1]{\tcbox[on line, boxrule=0pt, colback=accentsoft, colframe=accentsoft, arc=8pt,
  boxsep=0pt, left=4.5pt, right=4.5pt, top=3pt, bottom=3pt]{\makebox[0.62em]{\small\bfseries\color{accent}#1}}}
\newlist{etapes}{enumerate}{1}
\setlist[etapes]{label=\pastille{\arabic*}, leftmargin=28pt, labelsep=8pt, itemsep=5pt, topsep=4pt}

% Encadrés
\newtcolorbox{promesse}{enhanced, colback=oksoft, colframe=oksoft, arc=8pt, boxrule=0pt,
  left=12pt, right=12pt, top=9pt, bottom=9pt}
\newtcolorbox{note}{enhanced, colback=soft, colframe=soft, arc=8pt, boxrule=0pt,
  left=12pt, right=12pt, top=8pt, bottom=8pt, before skip=10pt, after skip=6pt}

% Règles : à faire / à ne pas faire
\newcommand{\regle}[4]{% couleur fond, couleur signe, signe, contenu
  \begin{tcolorbox}[enhanced, colback=white, colframe=line, boxrule=0.6pt, arc=8pt,
    left=40pt, right=12pt, top=8pt, bottom=8pt, before skip=0pt, after skip=7pt,
    overlay={\node[circle, fill=#1, text=#2, minimum size=18pt, inner sep=0pt, font=\bfseries]
      at ([xshift=20pt, yshift=-17pt]frame.north west) {#3};}]
  #4
  \end{tcolorbox}}
\newcommand{\afaire}[2]{\regle{oksoft}{ok}{{\symb ✓}}{\textbf{#1}\\[1pt]\muted{#2}}}
\newcommand{\pasfaire}[2]{\regle{nosoft}{no}{{\symb ✕}}{\textbf{#1}\\[1pt]\muted{#2}}}

% Écrans
\newcommand{\ecran}[2]{\begin{tcolorbox}[colback=white, colframe=line, boxrule=0.6pt, arc=8pt,
  left=10pt, right=10pt, top=7pt, bottom=7pt, before skip=0pt, after skip=0pt, height=17mm, valign=center]
  \textbf{#1}\\[1pt]{\small\muted{#2}}\end{tcolorbox}}

\begin{document}

% ---------- Titre ----------
\begin{tikzpicture}
  \node[fill=accent, rounded corners=6pt, minimum size=13mm, text=white, font=\displayfont\fontsize{22}{22}\selectfont] {N};
\end{tikzpicture}

\vspace{6pt}
{\displayfont\fontsize{30}{34}\selectfont Nini, mode d’emploi\par}
\vspace{4pt}
{\large\muted{Nini est ton carnet de travail : tes dossiers, les tâches du mois et ton planning, au même endroit.}\par}

\vspace{8pt}
\begin{promesse}
\textbf{Tout reste sur ton ordinateur.}\enspace Nini ne se connecte jamais à Internet. Aucune donnée ne sort.
\end{promesse}

% ---------- Écrans ----------
\section{Ce que tu trouves dans Nini}
\begin{tcbraster}[raster columns=2, raster equal height, raster column skip=8pt, raster row skip=8pt]
  \ecran{Aujourd’hui}{Les dossiers prévus ce jour, les retards, la semaine.}
  \ecran{Répartition}{Le planning du mois : quel dossier, quel jour, combien de temps.}
  \ecran{Toutes les tâches}{La liste complète, avec des filtres.}
  \ecran{Dossiers}{Chaque client, son manager, son temps alloué.}
  \ecran{Rapports}{Ce qui est fait et ce qui reste, à imprimer ou en Excel.}
  \ecran{Réglages}{Sauvegardes et routine du mois.}
\end{tcbraster}

% ---------- Le mois type ----------
\section{Le mois type}
\muted{Chaque dossier passe par ces quatre périodes. Nini crée les tâches et les étapes à cocher pour chacune.}

\vspace{4pt}
\begin{tikzpicture}[x=0.995\linewidth/31]
  \def\h{3.05}
  % fonds des périodes
  \fill[accentsoft, rounded corners=6pt] (0,0) rectangle (10,\h);
  \fill[white] (10,0) rectangle (17,\h);
  \fill[ambersoft] (17,0) rectangle (24,\h);
  \fill[white, rounded corners=6pt] (24,0) rectangle (31,\h);
  \draw[line, rounded corners=6pt, line width=0.6pt] (0,0) rectangle (31,\h);
  \foreach \x in {10,17,24} \draw[line, line width=0.6pt] (\x,0) -- (\x,\h);
  % textes
  \newcommand{\periode}[5]{%
    \node[anchor=north west, text width=#2*\linewidth/31-14pt, align=left, inner sep=0pt] at (#1+0.35,\h-0.18) {%
      {\ttfamily\small\color{ink2}#3}\\[2pt]\textbf{#4}\\[2pt]{\footnotesize\color{ink2}#5}};}
  \periode{0}{10}{1 → 10}{Début de mois}{Banque, recettes, lettrage, 471, balance. À finir avant le 17.}
  \periode{10}{7}{11 → 17}{Milieu}{Achats, paie, OD, états manquants.}
  \periode{17}{7}{18 → 24}{TVA}{Déclarations du 19, du 21 et du 24. Révision.}
  \periode{24}{7}{25 → fin}{Fin de mois}{Tableaux de bord, finalisation, point manager.}
  % semaine réservée
  \fill[pink, rounded corners=1.5pt] (4,-0.5) rectangle (10,-0.36);
  \node[anchor=west, font=\footnotesize, text=ink2] at (10.3,-0.43) {Du 5 au 10 : semaine réservée au dossier de la période réservée.};
\end{tikzpicture}

% ---------- Début de mois ----------
\section[Le 1er jour du mois]{Le 1\textsuperscript{er} jour du mois}
\begin{etapes}
  \item Ouvre \btn{Répartition}.
  \item Clique sur \btn{Créer les tâches du mois}.
  \item Clique sur \btn{Proposer une répartition}.
  \item Ajuste si besoin : fais glisser un bloc vers un autre jour, ou clique dessus pour changer la durée.
\end{etapes}

% ---------- Chaque jour ----------
\section{Chaque jour}
\begin{etapes}
  \item Ouvre Nini : \btn{Au programme aujourd’hui} te montre tes dossiers du jour.
  \item Clique sur un dossier : la tâche de la période s’ouvre.
  \item Coche les étapes au fur et à mesure. Tout s’enregistre tout seul.
  \item Note ton temps : bouton \btn{Chrono}, ou saisis « 1h30 ».
\end{etapes}
\begin{note}
Raccourcis : \kbd{N} crée une tâche, \kbd{Ctrl} + \kbd{K} retrouve n’importe quoi.
\end{note}

% ---------- Changer d'appareil ----------
\section{Changer d’appareil}
\muted{Par exemple pour passer de l’ordinateur à la tablette.}
\begin{etapes}
  \item Sur l’appareil où tu as travaillé : \btn{Exporter une sauvegarde}.
  \item Transporte le fichier : câble, clé USB ou lecteur du cabinet.
  \item Sur l’autre appareil : \btn{Réglages → Restaurer}, puis choisis le fichier.
\end{etapes}
\begin{note}
\textbf{Une seule règle :} ne travaille jamais sur deux appareils en même temps. La restauration remplace tout ce qu’il y avait avant.
\end{note}

% ---------- iPad ----------
\section{Sur l’iPad}
\muted{L’iPad n’ouvre pas le fichier Nini.html. On installe Nini comme une app, une seule fois.}
\begin{etapes}
  \item Avec Internet, une seule fois : ouvre l’adresse de Nini dans \btn{Safari}.
  \item Touche \btn{Partager → Sur l’écran d’accueil → Ajouter}. L’icône bleue « N » apparaît.
  \item Ouvre Nini \btn{depuis son icône}, puis \btn{Réglages → Restaurer} avec ta sauvegarde reçue par AirDrop.
  \item C’est tout : coupe le Wi-Fi ou passe en mode avion, Nini fonctionne.
\end{etapes}
\begin{note}
\textbf{Attention :} ne supprime pas l’icône de Nini. Cela effacerait tes données sur l’iPad.
Pour sauvegarder depuis l’iPad : \btn{Exporter une sauvegarde}, puis AirDrop ou « Enregistrer dans Fichiers ».
\end{note}

% ---------- Sécurité ----------
\section{Les règles de sécurité}
\afaire{Sauvegarde chaque vendredi.}{Réglages → Exporter une sauvegarde. Range le fichier sur le lecteur du cabinet. Nini te le rappelle après 7 jours.}
\afaire{Verrouille ton ordinateur quand tu t’absentes.}{Touches \kbd{Windows} + \kbd{L}. Nini n’a pas de mot de passe : il est protégé par ta session.}
\afaire{Ouvre toujours Nini de la même façon.}{Même fichier, même endroit, même navigateur (Edge ou Chrome).}
\pasfaire{Ne vide pas les données du navigateur.}{Cela efface Nini. Si c’est arrivé : Réglages → Restaurer ta dernière sauvegarde.}
\pasfaire{N’utilise pas la navigation privée.}{Tout serait perdu à la fermeture de la fenêtre.}
\pasfaire{N’envoie pas une sauvegarde par mail perso, WhatsApp ou cloud perso.}{Elle contient les noms des clients. C’est un document confidentiel du cabinet.}

% ---------- En cas de souci ----------
\section{En cas de souci}
{\renewcommand{\arraystretch}{1.55}
\arrayrulecolor{line}
\begin{tabularx}{\linewidth}{@{}>{\bfseries\raggedright\arraybackslash}p{0.38\linewidth}>{\raggedright\arraybackslash\leavevmode\color{ink2}}X@{}}
\hline
Je ne vois plus mes dossiers & Réglages → Restaurer, puis choisis ta dernière sauvegarde. \\ \hline
J’ai coché une tâche par erreur & Clique sur « Annuler » dans le message en bas de l’écran, ou décoche-la. \\ \hline
Le planning ne me convient pas & Déplace les blocs à la main, ou clique sur « Nouvelle proposition ». \\ \hline
Un dossier arrive ou s’arrête & Dossiers → Nouveau dossier, ou Archiver dans la fiche du dossier. \\ \hline
La date de TVA d’un dossier change & Fiche du dossier → Modifier → Date limite TVA (19, 21 ou 24). \\ \hline
\end{tabularx}}

\vfill
{\small\muted{Les dates proposées par Nini sont des repères. En cas de doute sur une échéance, la règle du cabinet prime.}}

\end{document}
